\chapter{waLBerla in a Nutshell}\label{r02:ch:walberla}

waLBerla (\emph{widely applicable lattice Boltzmann from Erlangen}) is developed at the
chair for system simulation of FAU Erlangen-Nürnberg~\cite{walberla,walberla-repo}. It
started as a lattice Boltzmann (LBM) code and grew into a general framework for
simulations on block-structured grids: multiphase and free-surface flows, phase-field
models, rigid-body and particle dynamics (the MESA-PD module) and fluid--particle
coupling. Its design goal is near-perfect scaling, from a laptop to hundreds of
thousands of cores~\cite{schornbaum-blockforest}. This chapter introduces the concepts a
developer needs to follow the rest of the report. It uses the lattice Boltzmann method
as the running example, because that is also the workload measured on the K1.

\section{The lattice Boltzmann method in one page}\label{r02:sec:lbm}

The LBM~\cite{krueger-lbm} simulates a fluid on a regular grid of cells. Instead of
velocity and pressure, each cell stores $q$ \term{particle distribution functions} (PDFs)
$f_i$, one for each discrete velocity $\mathbf{c}_i$ of a \term{stencil}. The stencils
are named D$d$Q$q$: D3Q19 is three-dimensional with 19 velocities (the cell itself, 6
face neighbours and 12 edge neighbours). One time step consists of two parts:
\begin{equation}
  f_i(\mathbf{x} + \mathbf{c}_i\,\Delta t,\; t + \Delta t)
  \;=\; f_i(\mathbf{x}, t) \;-\; \omega\,\bigl(f_i(\mathbf{x}, t) - f_i^{\text{eq}}(\rho, \mathbf{u})\bigr).
  \label{r02:eq:lbm}
\end{equation}
\begin{description}
  \item[Collision] (right-hand side) relaxes each PDF towards an equilibrium
    $f_i^{\text{eq}}$ that depends only on the cell's density $\rho = \sum_i f_i$ and
    velocity $\mathbf{u} = \frac{1}{\rho}\sum_i f_i\mathbf{c}_i$. With a single relaxation
    rate $\omega$ this is the SRT (BGK) collision operator; MRT, TRT, cumulant and
    central-moment operators use more rates. The collision is purely local and costs a
    few hundred floating-point operations per cell.
  \item[Streaming] (left-hand side) moves each post-collision PDF to the neighbour cell in
    its direction. It is pure data movement.
\end{description}
waLBerla's generated sweeps fuse both into one \term{stream-collide} sweep in the
\term{pull} pattern: each cell gathers $f_i(\mathbf{x} - \mathbf{c}_i)$ from its
neighbours, collides, and writes the 19 results into a second array, \code{f_tmp}. The
two arrays are swapped after the sweep. Per cell and time step the sweep reads 19 and
writes 19 values, plus the density and the three velocity components if they are
stored: 42 values, 336 bytes in double precision. With a few hundred floating-point
operations per 336 bytes, the LBM is usually limited by memory bandwidth, not by
arithmetic. That fact returns in \cref{r02:ch:results}.

Performance is reported in \term{MLUP/s}, million lattice-cell updates per second.

\section{The block-structured domain}\label{r02:sec:blockforest}

waLBerla does not store one big array. The simulation domain is split into
\term{blocks}, and every block holds a regular grid of cells of the same size
(\cref{r02:fig:blockforest}). The blocks are the unit of everything:
\begin{itemize}
  \item \textbf{Distribution.} Each MPI process owns one or more blocks. A uniform
        grid with one block per process is the simplest case, and the one the K1
        benchmark uses. For grid refinement, the blocks form an octree (a
        \term{forest} of octrees, hence \cls{BlockForest}), and load balancing moves
        blocks between processes.
  \item \textbf{Data.} Simulation data such as fields are \term{block data}: one
        instance per block, registered once and looked up with a handle.
  \item \textbf{Work.} A \term{sweep} is a function that processes one block. The
        time loop calls every sweep for every local block.
\end{itemize}
The class most applications use is \cls{StructuredBlockForest}, a block forest whose
blocks carry a Cartesian cell grid. \code{blockforest::createUniformBlockGrid(...)}
creates one from a domain bounding box, the number of blocks per direction and the
number of cells per block, together with periodicity flags.

\begin{figure}[htbp]
  \centering
  % Block-structured domain (left) and one block with ghost layer and halo exchange (right).
\begin{tikzpicture}[x=1cm,y=1cm,font=\footnotesize]
  % ---------------------------------------------------------------- left: domain
  \begin{scope}[shift={(0,0.3)}]
    \def\bs{1.05}
    \foreach \i/\j/\p/\c in {0/0/0/psblue, 1/0/0/psblue, 2/0/1/psteal, 3/0/1/psteal,
                             0/1/2/psorange, 1/1/2/psorange, 2/1/3/psgray, 3/1/3/psgray} {
      \fill[\c!18] (\i*\bs,\j*\bs) rectangle ++(\bs,\bs);
      \node[font=\scriptsize] at (\i*\bs+0.5*\bs,\j*\bs+0.5*\bs) {P\p};
    }
    \draw[psgray!80,step=\bs] (0,0) grid (4*\bs,2*\bs);
    \draw[black!75,line width=0.9pt] (0,0) rectangle (4*\bs,2*\bs);
    \draw[psorange!90!black,line width=1.4pt] (0,0) rectangle (\bs,\bs);
    \node[r02label,anchor=north] at (2*\bs,-0.12)
      {domain: $4\times2$ blocks on 4 processes;\\the orange block is shown on the right};
  \end{scope}

  % ---------------------------------------------------------------- right: one block
  \begin{scope}[shift={(5.7,0.0)}]
    \def\s{0.34}
    % our block: 8x8 cells = 6x6 interior + 1 ghost layer
    \fill[psgray!25] (0,0) rectangle (8*\s,8*\s);
    \fill[white] (\s,\s) rectangle (7*\s,7*\s);
    \draw[r02cell,step=\s] (0,0) grid (8*\s,8*\s);
    \draw[black!80,line width=0.9pt] (\s,\s) rectangle (7*\s,7*\s);
    \draw[psorange!90!black,line width=1.4pt] (0,0) rectangle (8*\s,8*\s);
    \node[r02label,anchor=north] at (4*\s,-0.12) {block of P0:\\$6\times6$ cells + ghost layer};

    % east neighbour (same process): 3 columns shown, its first interior column is the border slice
    \begin{scope}[shift={(8*\s+1.7,0)}]
      \fill[psgray!25] (0,0) rectangle (\s,8*\s);
      \fill[psgray!25] (0,0) rectangle (3*\s,\s);
      \fill[psgray!25] (0,7*\s) rectangle (3*\s,8*\s);
      \fill[psblue!35] (\s,\s) rectangle (2*\s,7*\s);
      \draw[r02cell,step=\s] (0,0) grid (3*\s,8*\s);
      \draw[black!80,line width=0.9pt] (\s,\s) -- (3*\s,\s) (\s,\s) -- (\s,7*\s) -- (3*\s,7*\s);
      \node[r02label,anchor=north] at (1.5*\s,-0.12) {east neighbour\\(P0)};
      \coordinate (eastslice) at (1.2*\s,4*\s);
    \end{scope}
    % north neighbour (other process): 3 rows shown, its first interior row is the border slice
    \begin{scope}[shift={(0,8*\s+1.0)}]
      \fill[psgray!25] (0,0) rectangle (8*\s,\s);
      \fill[psgray!25] (0,0) rectangle (\s,3*\s);
      \fill[psgray!25] (7*\s,0) rectangle (8*\s,3*\s);
      \fill[psorange!35] (\s,\s) rectangle (7*\s,2*\s);
      \draw[r02cell,step=\s] (0,0) grid (8*\s,3*\s);
      \draw[black!80,line width=0.9pt] (\s,3*\s) -- (\s,\s) -- (7*\s,\s) -- (7*\s,3*\s);
      \node[r02label,anchor=west] at (8*\s+0.15,1.5*\s) {north neighbour (P2)};
      \coordinate (northslice) at (4*\s,1.2*\s);
    \end{scope}

    % halo exchange arrows: border slice -> ghost layer
    \draw[r02arrow,psblue!80!black,thick] (eastslice) -- (7.5*\s,4*\s)
      node[pos=0.5,above,font=\scriptsize,text=psblue!80!black,align=center] {local\\copy};
    \draw[r02arrow,psorange!90!black,thick] (northslice) -- (4*\s,7.5*\s)
      node[pos=0.45,right,font=\scriptsize,text=psorange!90!black] {MPI message};
  \end{scope}
\end{tikzpicture}

  \caption[Block-structured domain, ghost layers and halo exchange]{Left: a 2D domain
    split into $4\times2$ blocks, distributed over four MPI processes (colours). Right:
    one block with its interior cells (white) and one ghost layer (grey). Before a
    sweep reads its neighbours, the halo exchange copies the border slice of each
    neighbouring block into the matching ghost layer: by a local copy if both blocks
    belong to the same process, by MPI otherwise.}
  \label{r02:fig:blockforest}
\end{figure}

\section{Fields}\label{r02:sec:fields}

A \term{field} is the per-block array of the simulation state: a 3D grid of cells with
$F$ values per cell (e.g.\ $F = 19$ PDFs, $F = 3$ velocity components, $F = 1$ density),
surrounded by $g$ \term{ghost layers}. A field of $N_x\times N_y\times N_z$ interior cells
therefore allocates $(N_x+2g)(N_y+2g)(N_z+2g)\cdot F$ values. Cell coordinates of the
interior start at 0, so ghost cells have negative indices or indices $\geq N$.

The memory \term{layout} decides which index varies fastest
(\cref{r02:fig:fzyx}):
\begin{description}
  \item[\code{fzyx}] (structure of arrays): $x$ is contiguous, then $y$, then $z$, and
    the value index $f$ is outermost. All PDFs of direction $i$ form one contiguous 3D
    array. This is the layout of all generated LBM code in waLBerla, and the one SIMD
    code needs: four neighbouring cells along $x$ are four consecutive doubles, exactly
    one 256-bit vector.
  \item[\code{zyxf}] (array of structures): all values of one cell are contiguous.
\end{description}
In the v8 core (\cref{r02:sec:v8}), an \code{fzyx} field additionally pads every
$x$-line so that each line starts at an aligned address. The alignment is chosen per
architecture in \file{include/walberla/v8/memory/MemoryTraits.hpp}: 32 bytes for AVX,
64 for AVX-512, 16 for NEON. On RISC-V there is no special case, so the fallback
\code{__BIGGEST_ALIGNMENT__} applies, which is 16 bytes with both board compilers. The
RVV code does not depend on it (\cref{r02:sec:lim-alignment}).

\begin{figure}[htbp]
  \centering
  % fzyx layout of one x-line of a PDF field and how the generated kernel covers it.
\begin{tikzpicture}[x=1cm,y=1cm,font=\footnotesize]
  \def\w{0.66}   % cell width
  \def\h{0.52}   % cell height
  % ---- line (y, z, i): ghost, 10 interior cells, ghost, padding
  \begin{scope}
    \fill[psgray!25] (0,0) rectangle (\w,\h);
    \foreach \x in {0,...,7} {
      \pgfmathsetmacro{\xx}{(\x+1)*\w}
      \fill[psblue!22] (\xx,0) rectangle ++(\w,\h);
      \node[font=\scriptsize] at (\xx+0.5*\w,0.5*\h) {\x};
    }
    \foreach \x in {8,9} {
      \pgfmathsetmacro{\xx}{(\x+1)*\w}
      \fill[psorange!30] (\xx,0) rectangle ++(\w,\h);
      \node[font=\scriptsize] at (\xx+0.5*\w,0.5*\h) {\x};
    }
    \fill[psgray!25] (11*\w,0) rectangle (12*\w,\h);
    \fill[pattern=north east lines,pattern color=psgray!70] (12*\w,0) rectangle (13*\w,\h);
    \draw[r02cell] (0,0) grid[xstep=\w,ystep=\h] (13*\w,\h);
    \draw[black!80,line width=0.9pt] (\w,0) rectangle (11*\w,\h);
    \node[font=\scriptsize,anchor=east] at (-0.1,0.5*\h) {$f_i(\cdot,y,z)$};
    \node[font=\scriptsize] at (0.5*\w,0.5*\h) {$-1$};
    \node[font=\scriptsize] at (11.5*\w,0.5*\h) {10};

    % vector chunks
    \draw[decorate,decoration={brace,amplitude=4pt,mirror},thick,psblue!80!black]
      (1*\w+0.04,-0.08) -- (5*\w-0.04,-0.08)
      node[midway,below=5pt,font=\scriptsize,text=psblue!80!black,align=center]
      {\texttt{vle64\_v\_f64m1}\\4 cells, one vector};
    \draw[decorate,decoration={brace,amplitude=4pt,mirror},thick,psblue!80!black]
      (5*\w+0.04,-0.08) -- (9*\w-0.04,-0.08)
      node[midway,below=5pt,font=\scriptsize,text=psblue!80!black,align=center]
      {next iteration\\\texttt{ctr\_0 += 4}};
    \draw[decorate,decoration={brace,amplitude=4pt,mirror},thick,psorange!90!black]
      (9*\w+0.04,-0.08) -- (11*\w-0.04,-0.08)
      node[midway,below=5pt,font=\scriptsize,text=psorange!90!black,align=center]
      {scalar\\remainder};
    \node[font=\scriptsize,anchor=north,align=center,text=psgray!90!black] at (12.5*\w,-0.08)
      {line\\padding};
    \node[font=\scriptsize,anchor=north,align=center,text=psgray!90!black] at (0.5*\w,-0.08)
      {ghost};
  \end{scope}

  % ---- next line (y+1) and next direction, as offsets
  \begin{scope}[shift={(0,1.35)}]
    \fill[psgray!25] (0,0) rectangle (\w,\h);
    \fill[psgray!25] (11*\w,0) rectangle (12*\w,\h);
    \fill[pattern=north east lines,pattern color=psgray!70] (12*\w,0) rectangle (13*\w,\h);
    \draw[r02cell] (0,0) grid[xstep=\w,ystep=\h] (13*\w,\h);
    \node[font=\scriptsize,anchor=east] at (-0.1,0.5*\h) {$f_i(\cdot,y{+}1,z)$};
  \end{scope}
  \draw[r02arrow] (1.5*\w,\h+0.02) -- (1.5*\w,1.33)
    node[midway,right,font=\scriptsize] {\texttt{+\,\_stride\_f\_1}};
  \node[font=\scriptsize,anchor=west,align=left,text=psgray!90!black] at (13*\w+0.25,1.35+0.5*\h)
    {next plane: \texttt{+\,\_stride\_f\_2}\\next direction $i{+}1$: \texttt{+\,\_stride\_f\_3}};
  \node[font=\scriptsize,anchor=west,align=left] at (13*\w+0.25,0.5*\h)
    {$x$ contiguous,\\stride 1 (assumed)};
\end{tikzpicture}

  \caption[The fzyx layout and vectorization along x]{The \texttt{fzyx} layout of a PDF
    field, shown for one $x$-line of one direction $i$: $x$ is the contiguous index, so
    a 256-bit vector holds 4 consecutive \texttt{double} cells (8 \texttt{float}s). The
    generated kernel runs the vector loop over the first $\lfloor N_x/4\rfloor\cdot 4$
    cells and a scalar remainder loop over the rest. The strides
    \texttt{\_stride\_f\_1}, \texttt{\_stride\_f\_2} and \texttt{\_stride\_f\_3} (in
    elements) lead to the next line, the next plane and the next direction. The stride
    along $x$ is assumed to be 1.}
  \label{r02:fig:fzyx}
\end{figure}

\section{Sweeps, halo exchange and the time loop}\label{r02:sec:sweeps}

A \term{sweep} is any callable that takes a block, \code{void operator()(IBlock*)}.
Typical sweeps are the stream-collide update, an initialization, a boundary condition
or the computation of an output quantity. Sweeps only update interior cells; they may
read the ghost layers.

The ghost layers are filled by \term{communication}, called \term{halo exchange} in the
v8 core. For every pair of neighbouring blocks and every direction, a \term{pack info}
packs the \term{border slice} of the sender (the outermost interior cells) into a buffer
and unpacks it into the \term{ghost slice} of the receiver. If both blocks live on the
same process, the pack info copies directly (\code{localCopy}). For the LBM pull scheme a
specialized pack info (\code{streamPullSync}) sends only the PDFs that actually stream
across the face, a fraction of all 19. Periodic boundaries are handled the same way:
the block at the right end of the domain is the left neighbour of the one at the left
end.

A time step is therefore: halo exchange, then sweeps on all local blocks. waLBerla's
\cls{SweepTimeloop} manages such sequences with timing and logging. A plain C++ loop
works as well, and is what the K1 benchmark uses (\cref{r02:lst:app-timestep}).

\section{The classic core and the v8 core}\label{r02:sec:v8}

waLBerla currently contains two generations of its core data structures:
\begin{description}[style=nextline]
  \item[The classic modules in \file{src/}] \file{core} (types, MPI, logging, math),
    \file{blockforest}, \file{field} (\cls{GhostLayerField}), \file{communication},
    \file{timeloop}, \file{lbm} (hand-written LBM, now ``legacy''),
    \file{lbm_generated}, \file{gpu}, \file{vtk}, \file{mesa_pd} and about twenty more.
    They are compiled into one static library per module (\code{walberla_core},
    \code{walberla_field}, \dots). Their field classes are separate for CPU and GPU.
  \item[The experimental \emph{v8} core in \file{include/walberla/v8/}] A header-only
    library (\code{walberla::v8}, enabled with \cmakeopt{WALBERLA_ENABLE_V8CORE}) that
    is portable between CPU and GPU by construction:
    \begin{itemize}
      \item \code{memory::Field<T, F, MemTag>}: a distributed field whose per-block
            storage is managed by a \cls{BufferSystem}. The \term{memory tag} says where
            the memory lives: \code{memtag::host} (aligned CPU memory),
            \code{device}, \code{unified}, \code{pinned}.
      \item \code{memory::FieldView}: the view of one block's portion, with
            \code{data()}, \code{shape()} and \code{strides()}.
      \item \term{Execution tags} (\code{exectag::Serial}, \code{OpenMP}, \code{GPU})
            and \code{sweep::forAllCells(xtag, interval, lambda)} for hand-written
            loops.
      \item \code{HaloExchange::create<Stencil, MemTag>(blocks).sync(...).build()}:
            the halo exchange, with generic pack infos for fields and for the LBM pull
            scheme.
    \end{itemize}
    The v8 core builds on the classic \file{core}, \file{blockforest}, \file{stencil}
    and \file{field} modules. The new examples in \file{apps/examples/} (e.g.\
    \code{DoubleShearLayer}, \code{FlowAroundSphere}) use it, and it is the field system
    the RVV benchmark uses.
\end{description}
Code generated by sweepgen (\cref{r02:ch:codegen}) can target either field system. A
generator script selects v8 fields with \code{Sweep.use_v8core_fields()}.

\section{Anatomy of a v8 application}\label{r02:sec:app-anatomy}

\Cref{r02:lst:app-setup,r02:lst:app-timestep} show the important parts of the new
benchmark application \file{apps/benchmarks/SweepgenLbmCPU/SweepgenLbmCPU.cpp}. It is
the \code{DoubleShearLayer} example reduced to what is timed. The classes in namespace
\code{gen} (\code{SetInitialState}, \code{LBM::InitPdfs}, \code{LBM::StreamCollide}) are
generated at build time from a Python script (\cref{r02:sec:sweepgen-script}).

\rneedspace{12\baselineskip}
\begin{ccode}[nolol=false,caption={[Setting up blocks, fields and generated sweeps]Setting up blocks, fields and generated sweeps (\texttt{SweepgenLbmCPU.cpp}, abridged).},label=r02:lst:app-setup]
using MemoryTag  = memtag::host;
using LbStencil  = stencil::D3Q19;
using PdfField_T = memory::Field< real_t, LbStencil::Q, MemoryTag >;   // 19 PDFs per cell

//  One block of cells^3 per MPI process, periodic in all directions
auto blocks = blockforest::createUniformBlockGrid(
   domainAabb, numBlocks[0], numBlocks[1], numBlocks[2],   // blocks per direction
   cells, cells, cells,                                    // cells per block
   /*oneBlockPerProcess*/ true, true, true, true);         // periodic in x, y, z
ScalarField_T rho{ *blocks };       // density
VectorField_T u{ *blocks };         // velocity, 3 components
PdfField_T    pdfs{ *blocks };      // fzyx layout, one ghost layer by default

//  Generated sweeps: plain C++ classes, constructed from fields and parameters
gen::SetInitialState setInitialState{ blocks, rho, u, delta, kappa, u_0 };
gen::LBM::InitPdfs   lbInit{ pdfs, rho, u };
for (auto& b : *blocks) { setInitialState(&b); lbInit(&b); }

gen::LBM::StreamCollide streamCollide{ pdfs, rho, u, omega };
auto haloExchange = HaloExchange::create< LbStencil, MemoryTag >(blocks)
                       .sync(halo_exchange::streamPullSync< LbStencil >(pdfs))
                       .build();
\end{ccode}

\rneedspace{12\baselineskip}
\begin{ccode}[nolol=false,caption={[The timed time step of the benchmark]The timed time step (\texttt{SweepgenLbmCPU.cpp}).},label=r02:lst:app-timestep]
haloExchange();   // the pull stream reads the ghost layers, so fill them first

auto timestep = [&]() {
   for (auto& b : *blocks)
      streamCollide(&b);   // generated: stream + collide on one block, then swap f <-> f_tmp
   haloExchange();
   sweep::sync();          // waits for GPU work; a no-op on the CPU
};
\end{ccode}

The type \code{real_t} is \code{double} or \code{float}, depending on the CMake option
\cmakeopt{WALBERLA_DOUBLE_ACCURACY}. The generated code follows the same switch
(\cref{r02:sec:buildconfig}), which is how the benchmark runs in both precisions.
